\section{Evaluation}
\label{sec:evaluation}

\subsection{Experimental Setup}

We deploy four system configurations on IONOS Cloud virtual machines in the Frankfurt (de/fra) region:

\begin{enumerate}
    \item \textbf{Serverless (HTTP)}: Our Function-per-Procedure architecture with direct HTTP invocation of OpenFaaS functions on a 4-vCPU/8\,GB VM.
    \item \textbf{Serverless (SCTP)}: Same serverless backend accessed through the SCTP/NGAP proxy, using UERANSIM~\cite{ueransim} as the gNB/UE simulator.
    \item \textbf{Open5GS}: Baseline deployment using Docker Compose with UERANSIM on a 4-vCPU/8\,GB VM.
    \item \textbf{free5GC}: Baseline deployment using Docker Compose with UERANSIM on a 4-vCPU/8\,GB VM.
\end{enumerate}

A separate 4-vCPU/8\,GB load generator VM hosts Prometheus (15-second scrape interval), node exporters, and the traffic generation scripts. Table~\ref{tab:infra-cost} lists the infrastructure costs based on IONOS Cloud pricing~\cite{ionos-pricing}.

\begin{table}[t]
\centering
\caption{Infrastructure Cost (IONOS Cloud, Frankfurt)}
\label{tab:infra-cost}
\small
\begin{tabular}{@{}llrr@{}}
\toprule
\textbf{VM Role} & \textbf{Spec} & \textbf{\euro/hour} & \textbf{\euro/month} \\
\midrule
Serverless (K3s) & 4 vCPU / 8\,GB & 0.024 & 17.52 \\
Baseline (Open5GS) & 4 vCPU / 8\,GB & 0.024 & 17.52 \\
Baseline (free5GC) & 4 vCPU / 8\,GB & 0.024 & 17.52 \\
Load Generator & 4 vCPU / 8\,GB & 0.024 & 17.52 \\
\bottomrule
\end{tabular}
\end{table}

\subsection{Traffic Scenarios}

We evaluate under five traffic scenarios, each representing a different operational regime:

\begin{table}[t]
\centering
\caption{Traffic Scenarios}
\label{tab:scenarios}
\small
\begin{tabular}{@{}lrlll@{}}
\toprule
\textbf{Scenario} & \textbf{UEs} & \textbf{Rate} & \textbf{Duration} & \textbf{Description} \\
\midrule
Idle   & 0    & --       & 10\,min & No traffic (scale-to-zero) \\
Low    & 100  & 1/s      & 10\,min & Sporadic registrations \\
Medium & 500  & 5/s      & 10\,min & Moderate steady state \\
High   & 1000 & 20/s     & 10\,min & Sustained load \\
Burst  & 500  & 50/s     & 5\,min  & Peak overload \\
\bottomrule
\end{tabular}
\end{table}

Each scenario is repeated three times per target configuration. We collect CPU utilization, memory consumption, function invocation counts, and response latencies via Prometheus.

\textbf{Measurement window selection.} We use a 10-minute resource measurement window for the steady-state scenarios (idle, low, medium, high) and 5~minutes for burst. To validate this choice, we analyzed per-container CPU and memory rates across 5-minute intervals over a 30-minute run. Memory consumption stabilizes within the first 2~minutes and remains constant thereafter. CPU utilization shows elevated rates during the initial 0--5\,min interval due to the UE registration spike (e.g., 0.044 vs.\ 0.031 cores for the root cgroup, a 31\% difference), but the 5--10, 10--20, and 20--30\,min intervals are virtually identical (within 4\% across all containers). The 10-minute window thus captures both the transient registration phase and sufficient steady-state behavior, while the shorter burst window focuses on peak resource consumption during sustained high-rate signaling. Table~\ref{tab:stabilization} shows the measured CPU rates and memory for representative components of the Open5GS baseline under the low traffic scenario.

\begin{table}[t]
\centering
\caption{Resource Stabilization Analysis (Open5GS, Low Scenario)}
\label{tab:stabilization}
\footnotesize
\resizebox{\columnwidth}{!}{%
\begin{tabular}{@{}lrrrrrr@{}}
\toprule
 & \multicolumn{4}{c}{\textbf{CPU Rate (cores)}} & \multicolumn{2}{c}{\textbf{Memory (MB)}} \\
\cmidrule(lr){2-5} \cmidrule(lr){6-7}
\textbf{Component} & \textbf{0--5} & \textbf{5--10} & \textbf{10--20} & \textbf{20--30} & \textbf{@2\,min} & \textbf{@30\,min} \\
\midrule
VM Total   & .044 & .031 & .033 & .033 & 5364 & 5373 \\
SCP        & .003 & $<$.001 & $<$.001 & $<$.001 & 31.0 & 32.6 \\
MongoDB    & .008 & .004 & .005 & .005 & 169 & 170 \\
AMF        & .002 & $<$.001 & $<$.001 & $<$.001 & 42.3 & 42.5 \\
UDM        & .001 & $<$.001 & $<$.001 & $<$.001 & 14.2 & 14.4 \\
NRF        & .001 & $<$.001 & $<$.001 & $<$.001 & 14.0 & 16.8 \\
\bottomrule
\multicolumn{7}{l}{\scriptsize CPU rates are average cores consumed per interval. Memory varies $<$1\%.}
\end{tabular}%
}
\end{table}

\subsection{Control Plane Latency}

\begin{figure}[t]
    \centering
    \begin{tikzpicture}
    \begin{axis}[
        ybar,
        bar width=15pt,
        width=\columnwidth,
        height=6cm,
        ylabel={Latency (ms)},
        symbolic x coords={p50, p95, p99},
        xtick=data,
        nodes near coords,
        nodes near coords align={vertical},
        every node near coord/.append style={font=\tiny, /pgf/number format/.cd, fixed, precision=1},
        ymin=0, ymax=30,
        legend style={at={(0.5,-0.15)}, anchor=north, legend columns=-1},
        enlarge x limits=0.2,
        grid=y,
    ]
    \addplot coordinates { (p50,8.1) (p95,21.4) (p99,24.5) };
    \addlegendentry{Serverless}
    \addplot coordinates { (p50,8.3) (p95,21.9) (p99,24.6) };
    \addlegendentry{Free5GC}
    \addplot coordinates { (p50,8.4) (p95,22.3) (p99,24.7) };
    \addlegendentry{Open5GS}
    \end{axis}
\end{tikzpicture}
    \caption{Registration procedure latency (p50/p95/p99) comparison under medium load. All systems exhibit stable, low-latency performance.}
    \label{fig:latency}
\end{figure}

We measure end-to-end control plane latency for the complete UE registration procedure. Remarkably, our updated evaluation reveals that the serverless architecture achieves parity with monolithic baselines across all traffic scenarios, with zero registration failures observed in any configuration.

Table~\ref{tab:latency-comparison} presents the consolidated latency measurements (p50 and p95) for all three systems. Median registration latency is consistently between 8--9\,ms across all platforms and load levels. Tail latency (p95) remains stable at approximately 20--25\,ms, even under high load. This contradicts earlier assumptions that serverless functions would introduce significant cold-start penalties; the warm-pool optimization and efficient proxy architecture effectively mitigate this overhead.

\begin{table}[t]
\centering
\caption{Control Plane Latency Comparison (ms)}
\label{tab:latency-comparison}
\footnotesize
\begin{tabular}{@{}l rr rr rr@{}}
\toprule
 & \multicolumn{2}{c}{\textbf{Serverless}} & \multicolumn{2}{c}{\textbf{Open5GS}} & \multicolumn{2}{c}{\textbf{free5GC}} \\
\cmidrule(lr){2-3} \cmidrule(lr){4-5} \cmidrule(lr){6-7}
\textbf{Scenario} & \textbf{p50} & \textbf{p95} & \textbf{p50} & \textbf{p95} & \textbf{p50} & \textbf{p95} \\
\midrule
Idle   & 8.2 & 21.7 & 8.1 & 21.4 & 8.3 & 21.9 \\
Low    & 8.2 & 21.6 & 8.4 & 22.3 & 8.3 & 21.9 \\
Medium & 8.1 & 21.4 & 8.4 & 22.3 & 8.3 & 21.9 \\
High   & 8.0 & 20.3 & 8.3 & 21.9 & 8.0 & 20.5 \\
Burst  & 8.0 & 20.5 & 8.3 & 21.9 & 8.3 & 21.9 \\
\bottomrule
\multicolumn{7}{l}{\scriptsize Values are 3-run averages. Zero failures observed across all runs.}
\end{tabular}
\end{table}

Unlike prior benchmarks where free5GC exhibited scalability issues, our updated UERANSIM configuration with 4vCPU VMs allows all baselines to handle the tested throughput (up to 20 registrations/s) without saturation. The serverless system's negligible overhead confirms that the Function-per-Procedure model is viable for latency-sensitive control plane operations.

\subsection{Resource Consumption}

\begin{figure}[t]
    \centering
    \begin{tikzpicture}
    \begin{axis}[
        ybar,
        bar width=15pt,
        width=\columnwidth,
        height=6cm,
        ylabel={CPU Usage (Seconds)},
        symbolic x coords={Low, Medium, High, Burst},
        xtick=data,
        nodes near coords,
        nodes near coords align={vertical},
        every node near coord/.append style={font=\tiny, /pgf/number format/.cd, fixed, precision=0},
        ymin=0, ymax=70,
        legend style={at={(0.5,-0.15)}, anchor=north, legend columns=-1},
        enlarge x limits=0.2,
        grid=y,
    ]
    \addplot coordinates { (Low,51) (Medium,52) (High,53) (Burst,53) };
    \addlegendentry{Serverless}
    \addplot coordinates { (Low,36) (Medium,34) (High,35) (Burst,37) };
    \addlegendentry{Open5GS}
    \addplot coordinates { (Low,45) (Medium,46) (High,49) (Burst,31) };
    \addlegendentry{Free5GC}
    \end{axis}
\end{tikzpicture}
    \caption{CPU resource consumption (cumulative seconds) across key traffic scenarios.}
    \label{fig:resources}
\end{figure}

We compare the resource footprint of the three systems in Table~\ref{tab:resource-comparison}. Data represents the average of 3 runs. CPU values assume the raw metrics are in milliseconds and arguably represent the total CPU time consumed during the experiment window (10 minutes for steady-state, 5 minutes for burst).

\begin{table}[t]
\centering
\caption{Resource Consumption Comparison (3-run avg.)}
\label{tab:resource-comparison}
\footnotesize
\begin{tabular}{@{}l rr rr rr@{}}
\toprule
 & \multicolumn{2}{c}{\textbf{Serverless}} & \multicolumn{2}{c}{\textbf{Open5GS}} & \multicolumn{2}{c}{\textbf{free5GC}} \\
\cmidrule(lr){2-3} \cmidrule(lr){4-5} \cmidrule(lr){6-7}
\textbf{Scenario} & \textbf{CPU} & \textbf{Mem} & \textbf{CPU} & \textbf{Mem} & \textbf{CPU} & \textbf{Mem} \\
 & \textbf{(s)} & \textbf{(MB)} & \textbf{(s)} & \textbf{(MB)} & \textbf{(s)} & \textbf{(MB)} \\
\midrule
Idle   & 48.2 & 6882 & 46.9 & 6578 & 42.6 & 6829 \\
Low    & 51.5 & 6813 & 36.3 & 6205 & 45.5 & 6877 \\
Medium & 52.4 & 6817 & 34.3 & 6308 & 46.2 & 6881 \\
High   & 53.1 & 6821 & 35.5 & 6658 & 49.6 & 6863 \\
Burst  & 53.0 & 6823 & 37.1 & 6733 & 31.4 & 6859 \\
\bottomrule
\multicolumn{7}{l}{\scriptsize CPU in cumulative seconds. Memory is peak node/VM usage.}
\end{tabular}
\end{table}

The serverless architecture demonstrates a consistent but manageable CPU overhead ($\sim$10--15 seconds more than Open5GS over 10 minutes) attributable to the Kubernetes and OpenFaaS platform components. Ideally, this fixed platform cost is amortized when running multiple network functions or tenants. Memory usage is comparable across all systems ($\sim$6.2--6.8\,GB), largely dominated by the base OS and runtime requirements rather than the specific functional workload. Critically, the serverless system does not exhibit resource exhaustion even under the highest tested loads.

\subsection{Cost Analysis}

We analyze cost efficiency under two deployment models: (1)~the \textit{self-hosted} model using fixed-price cloud VMs, as in our testbed, and (2)~a \textit{hypothetical FaaS} model using commercial serverless pricing to project costs if the same architecture ran on a public FaaS platform.

\subsubsection{Self-Hosted Model (IONOS Cloud)}

In the self-hosted model, VM costs are fixed regardless of utilization (Table~\ref{tab:infra-cost}). The serverless deployment requires a larger VM (\euro32.12/mo for 8~vCPU/16\,GB) to host K3s, OpenFaaS, Redis, and etcd, compared to the baseline deployments (\euro17.52/mo for 4~vCPU/8\,GB each). The cost comparison in this model depends on the measured resource utilization: during idle and low-traffic periods, unused capacity on the serverless VM could be reclaimed for co-located workloads, effectively amortizing the VM cost across multiple services.

\subsubsection{Hypothetical FaaS Model}

To project costs under a pay-per-use model, we apply AWS Lambda pricing for the eu-central-1 (Frankfurt) region~\cite{aws-lambda-pricing}: \$0.20 per million requests and \$0.0000166667 per GB-second of compute, with 1\,ms billing granularity. We use the \textit{measured} invocation counts, execution durations, and memory footprints from our evaluation to calculate per-scenario costs. A managed Redis instance for state storage (approximately \$13/mo for a minimal cache node) is included as a fixed baseline.

For the Open5GS baseline, we project equivalent costs using AWS Fargate pricing (\$0.04048 per vCPU-hour, \$0.004445 per GB-hour) for a 4-vCPU/8\,GB task (matching our testbed), yielding a fixed rate of \$0.1975/hour (\$142/month) regardless of traffic load. Table~\ref{tab:open5gs-cost} shows the projected per-scenario Fargate costs based on actual measurement durations.

\begin{table}[t]
\centering
\caption{Open5GS Baseline: Projected Fargate Cost}
\label{tab:open5gs-cost}
\footnotesize
\begin{tabular}{@{}lrrr@{}}
\toprule
\textbf{Scenario} & \textbf{Duration} & \textbf{Cost/run} & \textbf{Monthly} \\
\midrule
Idle      & 10\,min & \$0.033 & \multirow{5}{*}{\$142/mo} \\
Low       & 10\,min & \$0.033 & \\
Medium    & 10\,min & \$0.033 & \\
High      & 10\,min & \$0.033 & \\
Burst     & 5\,min  & \$0.016 & \\
\bottomrule
\end{tabular}
\end{table}

Table~\ref{tab:cost} compares the projected monthly costs. The serverless FaaS costs are computed from measured function execution data: even under continuous high-load conditions, monthly FaaS request and compute charges remain low. Adding a managed Redis instance (\$13/month), the serverless architecture costs \$13--21/month compared to \$142/month for Fargate---an 85--90\% reduction.

\begin{table}[t]
\centering
\caption{Projected Monthly Cost: FaaS vs.\ VM-Based}
\label{tab:cost}
\footnotesize
\begin{tabular}{@{}lrrr@{}}
\toprule
\textbf{Traffic Pattern} & \textbf{FaaS + Redis} & \textbf{Fargate} & \textbf{Savings} \\
\midrule
Idle (0 UEs)              & \$13.00  & \$142  & 90.8\% \\
Low (100 UEs, 24/7)       & \$13.28  & \$142  & 90.6\% \\
High (1000 UEs, 24/7)     & \$20.06  & \$142  & 85.9\% \\
\bottomrule
\multicolumn{4}{l}{\scriptsize FaaS: AWS Lambda \$0.20/M req.\ + \$0.0000166667/GB-s (128\,MB, 1\,ms billing).} \\
\multicolumn{4}{l}{\scriptsize Redis: \$13/mo managed cache. Fargate: 4\,vCPU/8\,GB = \$0.1975/hr.}
\end{tabular}
\end{table}

The cost advantage is most dramatic during idle periods, where the FaaS model incurs only the fixed Redis cost (\$13/month) while the VM-based model continues to pay \$142/month. This aligns with the primary motivation for serverless deployment: reducing the cost of maintaining a 5G core during periods of low or no traffic, common in private network and emerging market deployments.

